\section{Finite-rank operators and signed measures}\label{sec:operators}

Fix an integer $p\ge1$ and the block $(M_p,d_p,o)$ constructed in
\Cref{lem:block-geometry}, with its anchors $A_p$, label map $h$, and
parameters from \eqref{eq:parameters}. We now begin the proof of
\Cref{thm:blocks}. An operator that approximates the anchors will give a
signed measure at each graph vertex. The measure will have bounded total
variation and a large integral against a bump determined by that vertex's
label. These two estimates are the inputs for the graph argument.

\subsection{A common finite support}

Finite rank does not imply that the range is supported on finitely many
evaluations. The following perturbation supplies such a support while
retaining the bounds needed later.

\begin{lemma}\label{lem:finite-support}
Suppose that a bounded finite-rank linear operator
$T_0:\F(M_p)\to\F(M_p)$ satisfies $\norm{T_0}\le p$ and
\[
 \norm{T_0\delta_a-\delta_a}<\frac12\qquad(a\in A_p).
\]
Then there are a finite set $S\subseteq M_p$ containing $o$ and a bounded
finite-rank linear operator $T:\F(M_p)\to\F(M_p)$ such that
\begin{equation}\label{eq:operator-bounds}
 \begin{gathered}
 T\F(M_p)\subseteq\operatorname{span}\{\delta_j:j\in S\},\qquad
 \norm{T}\le2p,\\
 \norm{T\delta_a-\delta_a}<1\qquad(a\in A_p).
 \end{gathered}
\end{equation}
\end{lemma}

\begin{proof}
If $T_0=0$, take $T=0$ and $S=\{o\}$; the assumed anchor errors already
give the required bounds. Otherwise choose a basis $v_1,\ldots,v_r$ of
the range of $T_0$. The coordinate functionals on this finite-dimensional
range are bounded. Composing them with $T_0$ gives bounded functionals
$\phi_i$ on $\F(M_p)$ such that
\[
 T_0u=\sum_{i=1}^r\phi_i(u)v_i.
\]
Choose $0<\eta<\min\{p,1/(2B)\}$. Finite linear combinations of
evaluations are dense in $\F(M_p)$, so we can choose such combinations
$w_i$ with
\[
 \sum_{i=1}^r\norm{\phi_i}\norm{w_i-v_i}<\eta.
\]
The operator $Tu=\sum_i\phi_i(u)w_i$ then satisfies
$\norm{T-T_0}<\eta$. Let $S$ be the union of $\{o\}$ and the finite
supports of the chosen combinations $w_i$. Its range has the asserted
support, and $\norm{T}\le p+\eta<2p$. Since
$\norm{\delta_a}=d_p(a,o)\le B$, for every $a\in A_p$ we also have
\[
 \norm{T\delta_a-\delta_a}
 \le\norm{T_0\delta_a-\delta_a}+\norm{T-T_0}\norm{\delta_a}
 <\frac12+\eta B<1.\qedhere
\]
\end{proof}

\subsection{Measures detected by label bumps}

For a source vertex $x$, define a function on all of $M_p$ by
\begin{equation}\label{eq:bump}
 b_x(t)=L\left(1-\frac{\norm{h(x)-h(t)}_\infty}{K}\right)_+,
 \qquad s_+=\max\{s,0\}.
\end{equation}
The label map is $1$-Lipschitz by \Cref{lem:block-geometry}. Hence $b_x$
takes values in $[0,L]$, has Lipschitz constant at most $L/K\le1$, and
equals $L$ at the anchor $a(x)$, whose label is $h(x)$. The small
Lipschitz constant will let us compare $T\delta_x$ with the nearby
anchor evaluation even when $\norm{T}$ is as large as $2p$.

For a signed measure $\nu$ on a finite set $S$, write
\[
 |\nu|(E)=\sum_{j\in E}|\nu(\{j\})|\quad(E\subseteq S),
 \qquad \norm{\nu}_{\TV}=|\nu|(S).
\]

\begin{lemma}\label{lem:measure-bounds}
Let $S\subseteq M_p$ be finite and contain $o$, and let
$T:\F(M_p)\to\F(M_p)$ be a bounded linear operator satisfying
\eqref{eq:operator-bounds}. There is a family $(\mu_x)_{x\in M_p}$ of
signed measures on $S$, each of total mass one, such that
\begin{equation}\label{eq:affine-pairing}
 \int g\,d\mu_x=(T\delta_x)(g-g(o))+g(o)
\end{equation}
for every real Lipschitz function $g$ on $M_p$. They satisfy
\begin{equation}\label{eq:tv-bound}
 \norm{\mu_x}_{\TV}\le V=1+4pB\qquad(x\in M_p),
\end{equation}
and, for every source vertex $x$,
\begin{equation}\label{eq:bump-mass}
 \int b_x\,d\mu_x\ge\frac{3L}{4}.
\end{equation}
\end{lemma}

\begin{proof}
For each $x$, express
$T\delta_x=\sum_{j\in S\setminus\{o\}}c_x(j)\delta_j$ and set
\[
 \mu_x(\{j\})=c_x(j)\quad(j\ne o),\qquad
 \mu_x(\{o\})=1-\sum_{j\in S\setminus\{o\}}c_x(j).
\]
The empty sum is zero when $S=\{o\}$. These masses sum to one, and
\[
 \int g\,d\mu_x
 =\sum_{j\in S\setminus\{o\}}c_x(j)\bigl(g(j)-g(o)\bigr)+g(o),
\]
which is \eqref{eq:affine-pairing}.

To estimate the total variation, let $g$ take the sign of
$\mu_x(\{j\})$ at each $j\in S$ and be zero off $S$, using
$\operatorname{sign}(0)=0$. Its values lie in $[-1,1]$, so separation by one gives
$\Lipnorm{g}\le2$. Thus
\[
 \begin{split}
 \norm{\mu_x}_{\TV}
 &=\int g\,d\mu_x\\
 &=(T\delta_x)(g-g(o))+g(o)
 \le2\norm{T\delta_x}+1
 \le1+4pB.
 \end{split}
\]
The last inequality uses $\norm{T}\le2p$ and
$\norm{\delta_x}=d_p(x,o)\le B$. In particular, the estimate includes
the mass at the basepoint.

For a source vertex $x$, \Cref{lem:block-geometry} gives
$d_p(x,a(x))\le R$. The anchor bound in \eqref{eq:operator-bounds}
therefore implies
\[
 \begin{split}
 \norm{T\delta_x-\delta_{a(x)}}
 &\le\norm{T(\delta_x-\delta_{a(x)})}
       +\norm{T\delta_{a(x)}-\delta_{a(x)}}\\
 &\le2pR+1.
 \end{split}
\]
Since $b_x(a(x))=L$, the affine pairing gives
\[
 \begin{split}
 \int b_x\,d\mu_x-L
 &=(T\delta_x-\delta_{a(x)})(b_x-b_x(o))\\
 &\ge-\frac{L}{K}(2pR+1).
 \end{split}
\]
Here $(2pR+1)/K\le1/5<1/4$, because $K=10p(R+1)$ and $p\ge1$.
This proves \eqref{eq:bump-mass}.
\end{proof}

The measures may have coefficients of either sign. Their common total
mass means that $\mu_x-\mu_y$ annihilates constants, so a scalar test
can be centered at $o$ without changing its pairing with this difference.
Signed representations also occur in Lipschitz almost extension;
\cite[Proposition~6 and Remark~7]{Naor21} distinguishes them from the
restrictions imposed by probability measures. Here the total-variation
bound is what permits arbitrary signs in $T$.
